\section{最优停时与 Snell 包络：秘书问题}

\subsection{背景介绍}

前三章讨论了 OST 在给定鞅过程和一个固定的停时规则下的表现——第二章展示了 OST 完美成立的情况，第三章展示了如何构造辅助鞅来应用 OST，第四章展示了 OST 三个充分条件全部失效时的反例行为。这些都属于 OST 的\textbf{描述性}应用：给定鞅和停时，OST 告诉我们 $\mathbb{E}[X_\tau]$ 等于（或不等於）$\mathbb{E}[X_0]$。

然而，概率论中还有一类与此对偶的问题：如果可以\textbf{自己选择}停时以最大化某个期望收益，最优策略是什么？这就是最优停时（optimal stopping）问题。其中秘书问题是最优雅的经典例子：$N$ 个候选人依次面试，每次面试后必须立即决定录用此人并结束面试，或拒绝此人（不可回头）继续面试。目标为最大化录用到最优（排名第 1）候选人的概率。

该问题的最优策略具有简洁的形式：前 $r-1$ 个候选人仅用于建立基准，之后一旦遇到一个比之前所有人都好的候选人（称为``纪录保持者''，record），立即录用。当 $N \to \infty$ 时，最优 $r^* \approx N/e$，成功概率 $\approx 1/e \approx 36.8\%$。这被称为``37\% 法则''。

一个自然的问题是：能否用 OST 来研究这个最优停时问题？答案是不能——秘书问题的核心过程不是鞅。但这一``不能''本身构成了本章的研究动因。

\subsection{理论分析}

\subsubsection{``纪录指示过程''不是鞅}

令 $I_n = 1$ 表示第 $n$ 个候选人是前 $n$ 个中最好的（即是一个 record），$I_n = 0$ 否则。候选人的到达顺序是 $1, 2, \dots, N$ 的一个随机排列，等概率 $1/N!$。由对称性，第 $n$ 个候选人在前 $n$ 个中排第一的概率为 $1/n$，即 $\mathbb{P}(I_n = 1) = 1/n$。此外，不同位置上的纪录事件是独立的\footnote{这是随机排列中``纪录值''的经典结论：设 $\sigma$ 为 $\{1,\dots,N\}$ 的均匀随机排列，则 $\mathbf{1}_{\sigma(k) = \min_{i \leqslant k} \sigma(i)}$ 相互独立。}：
\begin{equation}
  \mathbb{P}(I_{n+1} = 1 \mid I_1, \dots, I_n) = \mathbb{P}(I_{n+1} = 1) = \frac{1}{n+1}.
\end{equation}
条件期望 $\mathbb{E}[I_{n+1} \mid \mathcal{F}_n] = 1/(n+1)$ 是常数（不依赖于 $\mathcal{F}_n$），而 $I_n$ 要么是 0 要么是 1，绝不等于常数 $1/(n+1)$，故 $\{I_n\}$ 不是鞅。更准确地说，$I_n$ 构成一个独立但非同分布的 Bernoulli 序列，其均值 $\mathbb{E}[I_n] = 1/n$ 随 $n$ 衰减。这意味着\textbf{信息是有价值的}：越早的候选人越不可能是 record（因为比较集合小），越晚的候选人越可能是 record（但此时后悔已经来不及了）。正是这种``信息衰减''赋予了最优停时以非平凡的结构——OST 告诉我们``如果是鞅，任何停时都改变不了期望''，秘书问题之所以有非零的优化空间（37\% 远高于随机选的 $1/N$），恰恰因为其核心过程逃脱了 OST 的约束。

\subsubsection{Snell 包络——最优停时的统一框架}

尽管 $\{I_n\}$ 不是鞅，最优停时问题仍然可以通过鞅论的延伸工具——\textbf{Snell 包络}\cite{williams1991probability}——来统一刻画。对于带有收益函数 $g$ 的最优停时问题 $\sup_\tau \mathbb{E}[g(X_\tau)]$，定义\textbf{值函数}
\begin{equation}\label{eq:snell_def}
  V_n = \operatorname*{ess\,sup}_{\tau \geqslant n} \mathbb{E}[g(X_\tau) \mid \mathcal{F}_n],
\end{equation}
其中 supremum 取遍所有满足 $\tau \geqslant n$ 的停时。$V_n$ 表示`在时刻 $n$ 尚未停止的前提下，从 $n$ 开始采用最优策略能获得的最大期望收益'。Snell 包络具有以下核心性质：
\begin{enumerate}[noitemsep]
  \item \textbf{上鞅性}：$\{V_n\}$ 是使得 $V_n \geqslant g(X_n)$ 成立的最小 $\{\mathcal{F}_n\}$-上鞅。这在直观上是合理的——等待的期望价值不可能一直增长，必须受到某种``贴现''。
  \item \textbf{最优停时的刻画}：最优停时就是 Snell 包络首次触及收益函数的时刻，
  \begin{equation}\label{eq:opt_stop}
    \tau^* = \inf\{n: V_n = g(X_n)\}.
  \end{equation}
  即在 $\tau^*$ 之前 $V_n > g(X_n)$（继续等待严格更好），在 $\tau^*$ 处 $V_n = g(X_n)$（立即停止与继续等待无差异）。
  \item \textbf{Doob 分解}：由 Doob-Meyer 定理，上鞅可唯一分解为 $V_n = M_n - A_n$，其中 $\{M_n\}$ 为鞅，$\{A_n\}$ 为可料增过程（$A_0 = 0$）。增过程 $A_n$ 恰好衡量了``等待的时间成本''——随着时间的流逝，剩余的时间越来越少，最优期望收益逐渐被蚕食。最优策略在 $A_n$ 的变化速率超过某个阈值之前停止。
\end{enumerate}

\subsubsection{秘书问题的 Snell 包络}

对于秘书问题，以下使用 $n=1,\dots,N$ 表示候选人的到达序号。若第 $n$ 个候选人是 record，则其成为全局最优的条件概率为
\begin{equation}
  g_n = \mathbb{P}(\text{第 } n \text{ 个候选人为全局最优}\mid \text{它是 record}) = \frac{n}{N};
\end{equation}
若它不是 record，则立即录用必然失败。精确的 Snell 包络可通过逆向递推计算：设 $V_n$ 为第 $n$ 个候选人是 record 时的最优条件成功概率。从 $n=N$ 向前递推：
\begin{align}
  V_N &= 1 \quad (\text{最后一个候选人若为 record，则必为全局最优}), \\
  V_n &= \max\Big\{g_n,\; \sum_{k=n+1}^{N}
  \mathbb{P}(\text{下一个 record 在 } k) V_k\Big\}.
\end{align}
下一个 record 出现在位置 $k$（$k > n$）的概率为 $\frac{n}{k(k-1)}$\footnote{推导：$k$ 是 record 且 $n+1,\dots,k-1$ 都不是 record。$\mathbb{P}(I_k = 1) \cdot \prod_{j=n+1}^{k-1} \mathbb{P}(I_j = 0) = 1/k \cdot \prod_{j=n+1}^{k-1} (1 - 1/j) = n/(k(k-1))$。}，故
\begin{equation}\label{eq:snell_rec}
  V_n = \max\Bigg\{\frac{n}{N},\; \sum_{k=n+1}^{N}
  \frac{n}{k(k-1)} V_k\Bigg\}.
\end{equation}
最优停止区域为使得 $\frac{n}{N} \geqslant \sum_{k=n+1}^{N} \frac{n}{k(k-1)} V_k$ 的第一个 $n$——即``立即录用''的收益第一次大于或等于``继续等待''的期望收益。可以解析证明，当 $N \to \infty$ 时，这一阈值趋近于 $\lceil N/e \rceil$。

\subsubsection{OST 为什么不冲突}

现在可以清晰地看出 OST 与秘书问题的关系。OST 说：如果一个\textbf{过程}是鞅，那么\textbf{任意}停时下 $\mathbb{E}[X_\tau] = \mathbb{E}[X_0]$。秘书问题中的过程不是鞅，所以 OST 没有做出任何断言。OST 的约束力体现在它的``反面''上——在对照组实验中，我们用零均值正态噪声（一个真正的鞅序列）取代相对排名，同样的``前 $r$ 期观察然后停止''策略下，成功率始终接近 $1/N$，与 $r$ 无关。这正是 OST 的威力：当过程确实是鞅时，任何停时规则都会被 OST``锁死''在零期望收益上。秘书问题之所以有价值，是因为其过程不在 OST 的管辖范围内。

\subsection{数值实验设计与参数设置}

实验设置 $N \in \{20, 50, 100, 200\}$，对每个 $N$ 扫描 $r/N \in (0, 1)$ 共 50 个取值点，每个点模拟 5000 次独立试验估计成功率。对于 Snell 包络可视化的实验，精确计算 $N=20$ 时的 $V_n$（通过式 \eqref{eq:snell_rec} 递推）。对照组实验用标准正态噪声序列替代相对排名，采用相同的策略观察成功率与 $r/N$ 的关系。

\subsection{结果与讨论}

\textbf{图~\ref{fig:ch05_success}} 展示了四个 $N$ 值下成功率随观察比例 $r/N$ 的变化曲线。每条曲线呈现鲜明的单峰形态：$r$ 太小时观察样本不足，基准太低，容易过早录用到不是全局最优的候选人（方差过大）；$r$ 太大时则剩余的选择太少，可能已经错过了最优候选人。最优 $r^*/N$ 出现在 $0.35$--$0.39$ 区间，随 $N$ 增大逐步趋近 $1/e \approx 0.368$。最大成功率也从 $N=20$ 时的 $\sim 0.39$ 逐渐接近 $1/e \approx 0.368$。Monte Carlo 散点与 $N \to \infty$ 时的理论极限曲线高度吻合。值得注意的一个细节是，$N=20$ 时的最优成功率略高于 $1/e$——这是因为有限 $N$ 下最优策略的收益通常略高于渐近极限（小 $N$ 时信息更集中）。

\textbf{图~\ref{fig:ch05_snell}} 展示了 $N=20$ 时精确计算的条件值函数 $V_n$（蓝色实线）与即时 payoff $g(n)$（红色虚线）。这里的 $V_n$ 表示“当前候选人已经是 record 时”的最优条件成功概率；因此它在阈值前由继续等待价值支配，约为 $0.384$，在阈值后与即时收益重合并随序号上升。两条曲线的首次交汇对应阈值 $r^*=7$：先观察前 7 位候选人，之后从第 8 位开始，一旦当前候选人是 record 就立即录用。$7/20 = 0.35$，接近 $1/e$。

\begin{figure}[ht]
  \centering
  \includegraphics[width=\linewidth]{ch05_success}
  \caption{秘书问题成功率与观察比例 $r/N$ 的关系。散点为 Monte Carlo 估计（每种 $N$ 各 5000 次试验）。实线为 $N \to \infty$ 时的理论曲线。灰色竖虚线标记最优观察比例 $r^*/N = 1/e$，灰色水平虚线标记最优成功率极限 $1/e$。}
  \label{fig:ch05_success}
\end{figure}

\begin{figure}[ht]
  \centering
  \includegraphics[width=\linewidth]{ch05_snell}
  \caption{Snell 包络条件值 $V_n$ 与即时 payoff $g(n)$ 的对比（$N=20$）。阈值前继续等待价值高于即时收益；阈值 $r^*=7$ 之后两者重合，表示一旦遇到 record（历史最优）就应立即停止。$7/20 = 0.35 \approx 1/e$。}
  \label{fig:ch05_snell}
\end{figure}

\textbf{图~\ref{fig:ch05_control}} 直接将秘书问题与鞅对照组并列对比，构成了本章的``核心信息图''。蓝色数据来自经典的秘书问题（$N=100$），呈现熟悉的单峰曲线，最优成功率接近 $1/e$。红色方格数据来自对照组——将候选人相对排名替换为零均值独立正态噪声，采用完全相同的策略（``前 $r$ 个观察，之后选第一个超过基准的''）。在鞅信号下，成功率恒为 $1/N = 0.01$（红色虚线），与 $r$ 完全无关。\textbf{这正是 OST 的``反面''体现}：当底层过程是鞅时，任何停时规则下期望收益等于初始期望——噪声序列的``全局最优''出现在每个位置的概率均等，所以无论用什么策略，选到全局最优的概率恒为 $1/N$。秘书问题之所以有意义，恰因为其过程不在此列。

\begin{figure}[ht]
  \centering
  \includegraphics[width=\linewidth]{ch05_control}
  \caption{秘书问题与鞅对照组的直接对比（$N=100$）。秘书问题（蓝色）：成功率呈现非平凡的 $r$ 依赖性，存在最优观察比例，最大成功率 $\approx 1/e$。鞅噪声对照组（红色）：成功率恒为 $1/N$（红色虚线），与 $r$ 完全无关——OST 的``反面''含义：当过程是鞅时，任何停时策略都无法改变期望。}
  \label{fig:ch05_control}
\end{figure}
